%==============================================================
%  Figura 6.3 - Cicli di vita delle macchine e delle SRP/CS
%  Adattamento in italiano della Figura 6.3 dell'IFA Report 2/2017
%
%  Le coordinate sono in decimi di millimetro con l'asse y rivolto
%  verso il basso (x=0.1mm, y=-0.1mm): corrispondono ai pixel di un
%  ritaglio dell'originale largo 1560 px, quindi per ritoccare una
%  posizione basta confrontarla con l'originale. La figura viene
%  stampata quasi in scala 1:1, per cui i corpi del testo sono quelli
%  effettivi (l'originale usa caratteri molto piccoli nei diagrammi).
%
%  \rbox{x1}{y1}{x2}{y2}{stile}{font}{testo}  riquadro con testo centrato
%  \dmd{cx}{cy}{larghezza}{altezza}{stile}{font}{testo} rombo
%
%  Nota: dopo \textsubscript il testo che segue "\\" va racchiuso
%  tra graffe, altrimenti la "(" viene letta da TikZ come coordinata.
%==============================================================
\begin{tikzpicture}[
  x=0.1mm, y=-0.1mm,
  font=\footnotesize,
  grigio/.style={draw=black, fill=black!13, line width=0.5pt},
  bianco/.style={draw=black, fill=white, line width=0.5pt},
  linea/.style={draw=black, line width=0.6pt},
  ar/.style={linea, -{Triangle[length=4pt,width=4pt]}},
  arv/.style={draw=black, line width=1.3pt, -{Triangle[length=4pt,width=4pt]}},
  arp/.style={draw=black, line width=1.1pt, densely dotted, -{Triangle[length=3.5pt,width=3.5pt]}},
  cornice/.style={draw=black, line width=0.9pt},
  lab/.style={inner sep=1pt}
]

% corpi del testo: fA etichette principali, fB flusso del Capitolo 6,
% fC modello a V e diagramma di validazione
\def\fA{\fontsize{8}{9.6}\selectfont}
\def\fB{\fontsize{5.6}{6.6}\selectfont}
\def\fC{\fontsize{5}{6}\selectfont}

\def\rbox#1#2#3#4#5#6#7{%
  \pgfmathsetmacro\tw{(#3-#1)*0.1-1.0}%
  \draw[#5] (#1,#2) rectangle (#3,#4);
  \node[align=center, text width=\tw mm, inner sep=0pt, font=#6]
    at ({(#1+#3)/2},{(#2+#4)/2}) {#7};}

\def\dmd#1#2#3#4#5#6#7{%
  \draw[#5] ({#1-#3/2},#2) -- (#1,{#2-#4/2}) -- ({#1+#3/2},#2) -- (#1,{#2+#4/2}) -- cycle;
  \node[align=center, inner sep=0pt, font=#6] at (#1,#2) {#7};}

%==============================================================
% Cornice generale
%==============================================================
\draw[line width=0.6pt] (5,5) rectangle (1555,1995);

%==============================================================
% Ciclo di vita della macchina secondo EN ISO 12100
%==============================================================
\draw[cornice] (73,40) rectangle (1537,585);

\rbox{114}{85}{733}{147}{grigio}{\fA}{Costruzione}
\rbox{114}{196}{733}{258}{grigio}{\fA}{Trasporto, montaggio e installazione}
\rbox{114}{307}{733}{369}{grigio}{\fA}{Messa in servizio}
\draw[ar] (424,147) -- (424,196);
\draw[ar] (424,258) -- (424,307);
\draw[linea] (424,369) -- (424,403) -- (806,403) -- (806,70) -- (1187,70);
\draw[ar] (1187,70) -- (1187,95);

\draw[grigio] (878,95) rectangle (1497,425);
\node[anchor=north west, align=left, inner sep=0pt, font=\fA] at (910,112)
  {Uso:\\[1pt]
   \renewcommand{\arraystretch}{1}%
   \begin{tabular}{@{}l@{\hspace{2mm}}>{\raggedright\arraybackslash}p{54mm}@{}}
     - & Attrezzaggio, apprendimento/programmazione e/o cambio di processo \\
     - & Funzionamento \\
     - & Pulizia \\
     - & Manutenzione \\
     - & Ricerca guasti
   \end{tabular}};
\draw[ar] (1187,425) -- (1187,462);
\rbox{878}{462}{1497}{560}{grigio}{\fA}{Messa fuori servizio, smantellamento e smaltimento}

\node[anchor=west, align=left, lab, font=\fA] at (95,530)
  {Ciclo di vita della macchina\\ secondo EN ISO 12100};

%==============================================================
% Collegamento e processo di riduzione del rischio
%==============================================================
\draw[black, fill=black!30, line width=0.6pt] (349,632) -- (428,587) -- (507,632) -- (428,677) -- cycle;
\draw[linea] (349,632) -- (507,632);

\rbox{212}{699}{646}{862}{grigio, line width=1.5pt}{\fontsize{9.5}{11.5}\selectfont}
  {Processo di\\ riduzione del rischio\\ secondo EN ISO 12100}

\node[anchor=north west, align=left, lab, font=\fA] at (440,866)
  {Misura di protezione\\ sotto forma di sistema\\ di comando};

%==============================================================
% Progettazione di sistemi di comando sicuri: Capitolo 6
%==============================================================
\draw[cornice] (73,992) rectangle (741,1907);

\draw[ar] (428,862) -- (428,1017);

\rbox{185}{1017}{670}{1046}{grigio}{\fB}{Identificazione delle funzioni di sicurezza (SF)}
\rbox{185}{1069}{670}{1098}{grigio}{\fB}{Specifica delle caratteristiche di ciascuna SF}
\rbox{185}{1121}{670}{1150}{grigio}{\fB}{Determinazione del PL richiesto ({PL\textsubscript{r}})}
\draw[bianco] (428,1172) circle (0.7mm);
\rbox{185}{1195}{670}{1224}{grigio}{\fB}{Realizzazione delle SF, identificazione degli SRP/CS}
\rbox{185}{1247}{670}{1306}{grigio}{\fB}
  {Determinazione del PL degli SRP/CS da categoria,\\ {\textit{MTTF}\textsubscript{D}}, {\textit{DC}\textsubscript{avg}}, CCF}
\rbox{185}{1306}{670}{1335}{grigio}{\fB}{Software e guasti sistematici}

\dmd{428}{1421}{364}{96}{grigio}{\fB}{Verifica:\\ PL $\geq$ {PL\textsubscript{r}}}
\dmd{428}{1561}{364}{96}{grigio}{\fB}{Validazione:\\ requisiti soddisfatti?}
\dmd{428}{1703}{364}{96}{grigio}{\fB}{Tutte le SF\\ analizzate?}

\draw[ar] (428,1046) -- (428,1069);
\draw[ar] (428,1098) -- (428,1121);
\draw[ar] (428,1150) -- (428,1165);
\draw[ar] (428,1179) -- (428,1195);
\draw[ar] (428,1224) -- (428,1247);
\draw[ar] (428,1335) -- (428,1373);
\draw[ar] (428,1469) -- (428,1513);
\draw[ar] (428,1609) -- (428,1655);

% ritorni "no"
\draw[linea] (610,1421) -- (684,1421) -- (684,1172);
\draw[linea] (610,1561) -- (702,1561) -- (702,1172);
\draw[ar] (702,1172) -- (435,1172);
\draw[linea] (610,1703) -- (722,1703) -- (722,1083);
\draw[ar] (722,1083) -- (670,1083);

\foreach \y in {1421,1561,1703} {
  \node[lab, anchor=south west, font=\fB] at (616,\y) {no};
}
\foreach \y in {1482,1624,1763} {
  \node[lab, anchor=east, font=\fB] at (422,\y) {sì};
}

% ritorno al processo di riduzione del rischio
\draw[linea] (428,1751) -- (428,1795) -- (20,1795) -- (20,781);
\draw[ar] (20,781) -- (212,781);

% "Per ciascuna SF" con graffa
\rbox{84}{1330}{156}{1410}{draw=none, fill=black!13}{\fC}{Per\\ ciascuna\\ SF}
\draw[linea, decorate, decoration={brace, amplitude=6pt}] (180,1622) -- (180,1113);

\node[anchor=west, align=left, lab, font=\fA] at (91,1862)
  {Progettazione di sistemi di\\ comando sicuri: Capitolo 6};

%==============================================================
% Modello a V del software: punto 6.3
%==============================================================
\draw[cornice] (796,648) rectangle (1537,1062);

\node[align=center, lab, font=\fC] at (885,705) {Specifica\\ delle funzioni\\ di sicurezza};
\draw[arv] (830,745) -- (955,745);

\rbox{955}{690}{1105}{790}{grigio, rounded corners=4pt}{\fC}{Specifica del\\ software legato\\ alla sicurezza}
\rbox{1300}{695}{1430}{785}{grigio, rounded corners=4pt}{\fC}{Validazione}
\draw[arv] (1430,740) -- (1500,740);
\node[align=center, lab, font=\fC] at (1478,712) {Software\\ validato};

% freccia cava "Validazione"
\draw[black, fill=black!6, line width=0.5pt]
  (1295,730) -- (1150,730) -- (1150,720) -- (1110,740) -- (1150,760) -- (1150,750) -- (1295,750) -- cycle;
\node[font=\fC, inner sep=0pt] at (1222,740) {Validazione};

\rbox{985}{815}{1125}{865}{grigio, rounded corners=4pt}{\fC}{Progettazione\\ del sistema}
\rbox{1260}{815}{1390}{865}{grigio, rounded corners=4pt}{\fC}{Test di\\ integrazione}
\rbox{1030}{895}{1170}{945}{grigio, rounded corners=4pt}{\fC}{Progettazione\\ dei moduli}
\rbox{1215}{895}{1335}{945}{grigio, rounded corners=4pt}{\fC}{Test dei\\ moduli}
\rbox{1120}{975}{1235}{1025}{grigio, rounded corners=4pt}{\fC}{Codifica}

% ramo discendente e ascendente della V
\draw[arv] (970,790) -- (970,840) -- (985,840);
\draw[arv] (1000,865) -- (1000,920) -- (1030,920);
\draw[arv] (1045,945) -- (1045,1000) -- (1120,1000);
\draw[arv] (1235,1000) -- (1262,1000) -- (1262,945);
\draw[arv] (1335,920) -- (1355,920) -- (1355,865);
\draw[arv] (1390,840) -- (1410,840) -- (1410,785);

% corrispondenze di verifica (puntinate)
\draw[arp] (1260,840) -- (1125,840);
\draw[arp] (1215,920) -- (1170,920);
\draw[arp] (1040,815) -- (1040,790);
\draw[arp] (1090,895) -- (1090,865);
\draw[arp] (1150,975) -- (1150,945);

\node[anchor=south east, align=right, lab, font=\fA] at (1525,1052)
  {Modello a V\\ del software:\\ punto 6.3};

%==============================================================
% Validazione: Capitolo 7
%==============================================================
\draw[cornice] (796,1101) rectangle (1537,1960);

\draw[bianco] (1273,1172) ellipse [x radius=7mm, y radius=2.2mm];
\node[font=\fC] at (1273,1172) {Inizio};

\rbox{807}{1222}{955}{1300}{bianco}{\fC}{Liste dei guasti\\ (punto 7.1.3 e\\ Allegato C)}
\rbox{977}{1222}{1124}{1300}{bianco}{\fC}{Considerazioni di\\ progettazione\\ (Capitolo 6)}
\rbox{1199}{1222}{1347}{1300}{bianco}{\fC}{Piano di\\ validazione\\ (punto 7.1.2)}
\rbox{1378}{1222}{1525}{1300}{bianco}{\fC}{Principi di\\ validazione\\ (punto 7.1.1)}
\rbox{977}{1316}{1124}{1372}{bianco}{\fC}{Documenti\\ (punto 7.1.4)}
\rbox{977}{1380}{1124}{1482}{bianco}{\fC}{Criteri per\\ l'esclusione\\ dei guasti\\ (Allegato C)}
\rbox{1199}{1339}{1347}{1406}{bianco, line width=1.3pt}{\fC}{Analisi\\ (punto 7.1.5)}
\rbox{1310}{1517}{1458}{1583}{bianco, line width=1.3pt}{\fC}{Prove\\ (punto 7.1.6)}
\rbox{1151}{1749}{1395}{1804}{bianco}{\fC}{Rapporto di validazione\\ (punto 7.1.7)}

\dmd{1273}{1476}{150}{88}{bianco}{\fC}{Analisi\\ sufficiente?}
\dmd{1385}{1662}{150}{86}{bianco}{\fC}{Prove\\ superate?}

\draw[bianco] (1273,1857) ellipse [x radius=7mm, y radius=2.2mm];
\node[font=\fC] at (1273,1857) {Fine};

% flusso principale
\draw[ar] (1273,1194) -- (1273,1222);
\draw[ar] (1378,1261) -- (1347,1261);
\draw[ar] (1273,1300) -- (1273,1339);
\draw[ar] (1051,1300) -- (1051,1318);
\draw[ar] (881,1300) -- (881,1427) -- (977,1427);
\draw[ar] (1124,1350) -- (1199,1350);
\draw[ar] (1147,1350) -- (1147,1245) -- (1199,1245);
\draw[ar] (1124,1392) -- (1199,1392);
\draw[ar] (1165,1392) -- (1165,1277) -- (1199,1277);
\fill (1147,1350) circle (0.5mm);
\fill (1165,1392) circle (0.5mm);
\draw[ar] (1273,1406) -- (1273,1432);
\draw[ar] (1273,1520) -- (1273,1749);
\draw[ar] (1348,1476) -- (1385,1476) -- (1385,1517);
\draw[ar] (1385,1583) -- (1385,1619);
\draw[ar] (1385,1705) -- (1385,1749);
\draw[ar] (1460,1662) -- (1505,1662) -- (1505,1373) -- (1347,1373);
\draw[ar] (1273,1804) -- (1273,1835);

\node[lab, font=\fC, anchor=east] at (1269,1532) {sì};
\node[lab, font=\fC, anchor=south west] at (1352,1476) {no};
\node[lab, font=\fC, anchor=east] at (1381,1717) {sì};
\node[lab, font=\fC, anchor=south west] at (1464,1662) {no};

% riquadro degli aspetti da validare
\draw[bianco] (900,1500) rectangle (1112,1910);
\rbox{908}{1508}{1104}{1558}{grigio}{\fC}{Funzioni di sicurezza\\ (punto 7.3)}
\draw[grigio] (908,1566) rectangle (1104,1804);
\node[align=center, font=\fC, inner sep=0pt, anchor=north] at (1006,1576)
  {Performance Level (PL)\\ (punto 7.4)};
\node[anchor=north west, align=left, font=\fC, inner sep=0pt] at (922,1628)
  {- Categoria\\ - {\textit{MTTF}\textsubscript{D}}\\ - \textit{DC}\\ - CCF\\ - guasti sistematici\\ - Software};
\rbox{908}{1812}{1104}{1902}{grigio}{\fC}{Combinazione/\\ integrazione\\ (punto 7.6)}

% frecce tratteggiate verso il riquadro
\draw[linea, dashed, -{Triangle[length=4pt,width=4pt]}] (1197,1410) -- (1118,1494);
\draw[linea, dashed, -{Triangle[length=4pt,width=4pt]}] (1295,1570) -- (1116,1570);

\node[anchor=west, lab, font=\fA] at (816,1935) {Validazione: Capitolo 7};

%==============================================================
% Frecce curve grigie: rimando al modello a V e alla validazione
%==============================================================
% prima il contorno nero, poi il riempimento grigio sullo stesso percorso
\tikzset{
  curvabordo/.style={black, line width=6.4pt, -{Triangle[length=15pt,width=19pt]}},
  curvagrigia/.style={black!45, line width=5pt, shorten >=1.2pt,
                      -{Triangle[length=13pt,width=16.5pt]}}
}
\draw[curvabordo]  (600,1322) .. controls (760,1330) and (790,1120) .. (950,915);
\draw[curvagrigia] (600,1322) .. controls (760,1330) and (790,1120) .. (950,915);
\draw[curvabordo]  (555,1562) .. controls (555,1440) and (760,1420) .. (890,1490);
\draw[curvagrigia] (555,1562) .. controls (555,1440) and (760,1420) .. (890,1490);

\end{tikzpicture}
